\section{Evaluating use, retention and adaptation}
\label{sec:capability_eval}
The current development study measures autonomous tool use and portfolio outcomes.
The broader framework requires comparisons at three stages. The following design
specifies those comparisons; it does not report completed memory or context ablations.

\paragraph{Method selection and task quality.}
Tasks should expose method requirements and allow an agent to choose tools, parameters
and stopping points. A method can be appropriate without being uniquely optimal.
Selection can be assessed against task-specific eligibility and suitability criteria
defined independently of the agent output, alongside the quality of the resulting
forecast, allocation or analysis. Outcome-only rankings risk rewarding a lucky choice.
Matched controls distinguish basic inputs, skill guidance, fixed tool outputs and
agent-selected calls. A competing tool-access baseline is needed to isolate the
library's contribution from access to executable methods in general.

\paragraph{Context retention.}
With the backbone, available information and tools fixed, compare full context,
simple budget-matched retention and HiSTrim. Measure task quality, retention of
required evidence, allocated KV memory and wall-clock latency together. A smaller
context or higher confidence alone is not an improvement. This comparison also
requires defining financial record spans and protected task information before testing.

\paragraph{Experience and adaptation.}
With context and tool access fixed, compare no persistent memory, frozen memory,
an updating non-biological baseline and the proposed memory component. Matched
controllers should receive the same eligible feedback and action opportunities.
Track how choices change after feedback, performance on later unexamined tasks,
adaptation after shifts, turnover and cost. A separate feedback-check ablation tests
update eligibility; a matched exposure control separates learning from a risk gate.
Context-by-memory comparisons are necessary to assess their interaction rather than
attribute separate gains to the full system.

The primary measures should reflect the task: forecast accuracy or calibration,
decision utility, or the supported quality of an analysis. Report valid completion,
ineligible information use and execution failures as constraints on interpretation.
Freeze outputs before independent scoring, preserve failed attempts, and use new tasks
or chronological periods for evaluation. Report actual calls, tokens and latency,
and compare matched budgets; portfolio claims additionally need cost and risk controls.
